For a disclosure-controlled observed count, the end-to-end result envelope is
\begin{equation}
\mathcal{R}^{\mathrm{DC}}
=\left\{R(\mathbf O^{L}),R(\mathbf O^{*}),R(\mathbf O^{U})\right\},
\label{eq:confidentiality-envelope}
\end{equation}
where the three runs differ only in the declared count interpretation.

For uncertain parameter $j$ in a Latin-hypercube design with $M$ draws, let
$\pi_{mj}$ be a seeded random permutation of strata and $u_{mj}\sim U(0,1)$.
The unit-interval design and transformed parameter draw are
\begin{align}
q_{mj}
 &=\frac{\pi_{mj}+u_{mj}}{M},
\label{eq:lhs-stratum}\\
\psi_{mj}
 &=F_j^{-1}(q_{mj}),
\label{eq:lhs-transform}\\
\widehat{\mathbf y}^{(m)}
 &=\mathcal{M}\!\left(G,S,\boldsymbol\psi_m\right).
\label{eq:parameter-ensemble}
\end{align}
The marginal distributions or bounded ranges $F_j$, dependencies, network $G$,
and scenario $S$ are declared. Ensemble quantiles are not universal confidence
limits unless the parameter distributions have empirical probability meaning.

Across $M$ declared parameter or structural sensitivity runs, candidate $C$'s
top-$k$ rank acceptability is
\begin{equation}
\pi_C^{(k)}
=\frac{1}{M}\sum_{m=1}^{M}I\!\left(\operatorname{rank}_m(C)\le k\right).
\label{eq:rank-acceptability}
\end{equation}
